\chapter{Crypto-Native Firms}\label{ch:in:crypto-native-firms}

The largest US-listed crypto exchange reported 4\,510 employees at the end of 2022, after a restructuring in June that
had cut about 18\% of its staff; in January 2023 it cut about 21\% more, and a year later it had 3\,416. Its revenue had
fallen by 59\% in 2022. In 2024 its revenue more than doubled, and by the end of 2025 it employed 4\,951 people. In crypto,
headcount follows revenue, which follows the market for the assets the industry trades, with a lag of about a year;
and many employees are paid in those assets as well. This chapter describes crypto-native firms as employers: what kinds there are, what is different about
working in them, how their headcount moves with the cycle, what pay in tokens is worth, and what happens to an employee
when such a firm fails.

\section{Kinds of employer}

Crypto-native firms are those whose business began in crypto assets rather than moving into them.
\begin{itemize}
\item \textbf{Market makers and trading firms} quote on centralised exchanges and trade on chain (Book~11,
chapters~24--25), and make markets for token issuers under agreements (Book~3, chapter~24).
\item \textbf{Funds} invest in tokens and in the companies and protocols that issue them.
\item \textbf{Exchanges and brokers} run centralised venues (Book~3, chapter~15), custody and on-ramps between money and
crypto assets.
\item \textbf{Protocol teams and foundations} build and fund the software of a blockchain or an application on it.
\end{itemize}

\begin{definition}[Protocol foundation]\label{def:in:crypto-native-firms:foundation}
A \emph{protocol foundation}\index{protocol foundation} is a non-profit organisation that funds and coordinates the
development of an open-source blockchain protocol and its ecosystem, typically from a treasury held partly in the
protocol's own token, and that does not own the protocol it supports.
\end{definition}

One foundation describes itself as ``a non-profit that supports the Ethereum ecosystem'' and adds that it is ``not a tech
company, or a normal non-profit''. Its treasury policy shows what funding from a token means for the people it pays:
``periodically selling ETH to ensure sufficient savings for future years, and programmatically increasing our fiat
savings in bull markets to fund spending in bear markets''. A foundation's budget, and so its hiring, depends on the
price of what it holds.

\section{What is different: round-the-clock operations, custody, token pay, regulatory status}

Four things set crypto-native employers apart from the rest of the industry.
\begin{enumerate}
\item \textbf{The markets never close} (Book~3, chapter~29, calls the practice follow-the-sun trading): on-call rotas,
weekend incidents and global teams are the norm, not the exception (chapter~26).
\item \textbf{Custody is the firm's own problem.} Keys, wallets and the operational security around them are a core
engineering job (Book~3, chapter~24), and a failure is a loss, not an outage.
\item \textbf{Pay is often partly in tokens}, whose value moves with the market and whose tax treatment differs by
country (section~4).
\item \textbf{Regulatory status is still being settled} in many jurisdictions (Book~3, chapter~24, on the EU's
Markets in Crypto-Assets Regulation): firms move, license, and restructure as rules arrive, and employees move with them.
\end{enumerate}

\section{Headcount through the cycle}

\begin{dated}{2025-12}{One listed crypto exchange through the cycle}\label{dat:in:crypto-native-firms:coinbase}
Coinbase (Forms 10-K for 2022 to 2025): total revenue \$7\,839 million (2021), \$3\,194 million (2022), \$3\,108 million
(2023), \$6\,564 million (2024), \$7\,181 million (2025); employees at year end 4\,510 (2022), 3\,416 (2023), 3\,772 (2024),
4\,951 (2025). In June 2022 a restructuring affected about 18\% of headcount (about 1\,100 employees); in January 2023 a
second affected about 21\% of the end-2022 headcount (about 950 employees).
\end{dated}

The pattern is lagged (\cref{fig:in:crypto-native-firms:cycle}). Revenue fell 59\% in 2022 and 3\% in 2023; headcount
fell 24\% in 2023, the year after the revenue collapse. Revenue rose 111\% in 2024 and 9\% in 2025; headcount rose 10\% in
2024 and 31\% in 2025. A firm hires on last year's revenue and cuts on it too, so an employee joining at the top of a cycle
joins the cohort most exposed to the next cut. Five years of one firm is not a law, but it is the pattern a crypto job
seeker should expect until the evidence says otherwise.

\begin{omfigure}
\begin{tikzpicture}
\begin{axis}[width=10.5cm,height=5.4cm,xmin=2020.6,xmax=2025.4,ymin=0,ymax=9,xtick={2021,2022,2023,2024,2025},
  xticklabel style={/pgf/number format/1000 sep={}},ylabel={\small \$ billion; employees (thousands)},
  tick label style={font=\footnotesize},grid=major,grid style={gray!20},table/col sep=comma,unbounded coords=jump,
  legend style={legend cell align=left,font=\scriptsize,at={(0.5,-0.22)},anchor=north,/tikz/every even column/.append style={column sep=8pt}},legend columns=2]
\addplot[omDef,thick,mark=*] table[x=year,y=revenue]{figdata/industry/10-crypto-native-firms/coinbase.csv};
\addplot[omThm,thick,mark=square*] table[x=year,y=employees]{figdata/industry/10-crypto-native-firms/coinbase.csv};
\legend{total revenue (\$ billion),employees at year end (thousands)}
\end{axis}
\end{tikzpicture}

\omcaption{One listed crypto exchange, 2021--2025: total revenue and year-end employees. Employees for 2021 are not in the
filings read. Headcount follows revenue with about a year's lag. Data: Coinbase Forms 10-K for 2022--2025, through
\texttt{in\_crypto.headcount}.}\label{fig:in:crypto-native-firms:cycle}
\end{omfigure}

Revenue per employee shows the same cycle from the employee's side (\cref{fig:in:crypto-native-firms:perhead}): \$0.71
million per year-end employee in 2022, \$0.91 million in 2023, \$1.74 million in 2024 and \$1.45 million in 2025, a
swing by a factor of 2.5 in two years. Chapter~12 measured the same sensitivity across firms: crypto revenue per head
moved against the equity market's volatility, not with it. A bonus pool set from revenue moves with this line; a
headcount set from last year's revenue lags it, which is when cuts come.

\begin{omfigure}
\begin{tikzpicture}
\begin{axis}[ybar,width=10.5cm,height=4.8cm,xmin=2021.4,xmax=2025.6,ymin=0,ymax=2,bar width=16pt,
  xtick={2022,2023,2024,2025},xticklabel style={/pgf/number format/1000 sep={}},
  ylabel={\small \$ million per employee},tick label style={font=\footnotesize},grid=major,grid style={gray!20},
  table/col sep=comma,nodes near coords,nodes near coords style={font=\scriptsize},
  nodes near coords={\pgfmathprintnumber[fixed,precision=2]{\pgfplotspointmeta}}]
\addplot[fill=omDef,draw=none] table[x=year,y=perhead]{figdata/industry/10-crypto-native-firms/perhead.csv};
\end{axis}
\end{tikzpicture}

\omcaption{One listed crypto exchange's total revenue per year-end employee, 2022--2025. Data: Coinbase Forms 10-K,
through \texttt{in\_crypto.per\_head}.}\label{fig:in:crypto-native-firms:perhead}
\end{omfigure}

\section{Paid in tokens}

\begin{definition}[Token compensation, token warrant]\label{def:in:crypto-native-firms:tokens}
\emph{Token compensation}\index{token compensation} is pay delivered in a crypto asset, usually under a vesting schedule
with a cliff, and often with a lock-up after vesting during which vested tokens may not be sold. A \emph{token
warrant}\index{token warrant} is a right, granted with equity in a company, to receive tokens that the company or an
affiliate may issue in the future, in proportion to the equity held.
\end{definition}

Tax authorities treat tokens received for work as pay. The US tax authority's guidance states that virtual currency paid
by an employer as remuneration for services constitutes wages for employment tax purposes; the UK's states that
cryptoassets received as employment income count as money's worth and are subject to income tax and National Insurance on
their value. The value taxed is the value at receipt, usually at vesting. If the token then falls during a lock-up, the
employee has paid tax, in cash, on value that no longer exists.

\begin{method}[Valuing a token grant]\label{met:in:crypto-native-firms:grant}
\begin{enumerate}
\item Write the schedule: the cliff tranche, the monthly tranches, the lock-up after each vesting.
\item Simulate the token's price in the grant's currency over the vesting period and the lock-up, with a stated
volatility and drift (parameters, not forecasts).
\item For each tranche, compute the tax due at vesting and the proceeds at the first date it may be sold, net of any
liquidity discount.
\item Report the distribution of net proceeds against the same grant value paid in cash, and how often a tranche's tax
exceeds its proceeds.
\end{enumerate}
\end{method}

\begin{example}[A four-year grant]\label{ex:in:crypto-native-firms:grant}
A grant worth \$400\,000 at the token's price on the grant date vests 25\% after a one-year cliff and then monthly over
three years; each tranche is locked up for six months; tax is 40\% of its value at vesting (illustrative). With a
volatility of 80\% a year and no drift, the median net outcome over 20\,000 simulated paths is \$106\,193, against \$240\,000
for the same grant paid in cash; the 10th and 90th percentiles are \$17\,467 and \$537\,810; the mean, \$235\,062, is close
to the cash figure. In 76.2\% of paths at least one tranche's tax exceeds what it later sells for
(\cref{fig:in:crypto-native-firms:grant}). At a volatility of 40\% the median is \$194\,318 and such a tranche occurs in
2.4\% of paths.
\end{example}

\begin{omfigure}
\begin{tikzpicture}
\begin{axis}[width=10.5cm,height=5.4cm,ybar,bar width=4.5pt,xmin=-100,xmax=1000,ymin=0,
  xlabel={\small net outcome of the grant (\$ thousand)},ylabel={\small share of paths (\%)},
  tick label style={font=\footnotesize},grid=major,grid style={gray!20},table/col sep=comma]
\addplot[fill=omDef!55,draw=omDef] table[x=mid,y=share]{figdata/industry/10-crypto-native-firms/grant_hist.csv};
\draw[omThm,thick,dashed] (axis cs:240,0) -- (axis cs:240,26);
\node[font=\scriptsize,omThm,anchor=south west] at (axis cs:245,20) {cash equivalent \$240k};
\end{axis}
\end{tikzpicture}

\omcaption{The net outcome of a \$400\,000 four-year token grant (25\% cliff at one year, monthly thereafter, six-month
lock-up, 40\% tax at vesting), 20\,000 simulated price paths at 80\% annual volatility and no drift; outcomes above
\$1 million are counted in the last bin. The dashed line is the same grant paid in cash. Illustrative parameters. Data:
\texttt{firm.tokencomp.simulate}, through \texttt{in\_crypto.grant\_hist}.}\label{fig:in:crypto-native-firms:grant}
\end{omfigure}

The distribution is skewed: most paths leave the employee with less than cash would have, a minority with much more.
The mean is close to the cash value because the price has no drift, which is exactly why the median is not: at 80\%
volatility the typical path loses value. A candidate comparing a token grant with cash is comparing a lottery of about
equal expectation with a certainty, before counting the tax that falls on value at vesting.

\section{When the firm fails: the employee in the bankruptcy record}

Crypto firms can fail quickly (Book~3, chapter~15, and Book~16, chapter~29, describe a large exchange failure). For an
employee, a failure turns pay into a claim. In a US bankruptcy,
unpaid wages, salaries and commissions earned in the 180 days before the petition have priority among unsecured claims
only up to a fixed amount per individual, \$17\,150 since 1 April 2025; the rest ranks with other unsecured creditors.
Unvested tokens and equity in the failed firm are usually worth nothing. The practical lessons are those of any young
firm, sharper: take cash where possible, sell vested tokens on the schedule you planned, and know which entity employs
you.

\section{Tutorial: headcount and a token grant}\label{tut:in:crypto-native-firms:tut}

\begin{tutorial}
\textbf{Goal.} Read one crypto exchange's headcount against its revenue, and value a token grant.
\textbf{End state:} \cref{fig:in:crypto-native-firms:cycle,fig:in:crypto-native-firms:grant}.
\begin{tutsteps}
\item \textbf{The filings.} \texttt{data/industry/crypto\_headcount.csv} holds year-end employees and total revenue from
the exchange's Forms 10-K; \texttt{in\_crypto.changes()} computes the yearly changes.
\item \textbf{The grant.} \texttt{firm.tokencomp.Grant(400\_000, months=48, cliff=12, lockup=6, tax=0.40)} and
\texttt{simulate(grant, vol, drift, n, rng)} (\cref{lst:in:crypto-native-firms:sim}).
\omcode{code/firm/tokencomp/firm_tokencomp.py}{39}{58}{Each tranche taxed at vesting and sold after its lock-up, over
simulated price paths.}\label{lst:in:crypto-native-firms:sim}
\item \textbf{Summarise.} \texttt{firm.tokencomp.summary} gives the median, the 10th and 90th percentiles, the mean,
the cash equivalent and the share of paths with a tranche whose tax exceeds its proceeds.
\end{tutsteps}
\end{tutorial}

\textbf{What to change next.} Lengthen the lock-up to twelve months and see how often tax exceeds proceeds
(exercise~7); give the token a positive drift and ask what drift makes the median equal to the cash value.

\section{Build: the token-grant model}\label{bld:in:crypto-native-firms:tokencomp}

\begin{build}
\bfield{Purpose} Put a number, and a distribution, on pay in tokens, for this chapter and for chapter~13's comparison of
pay packages.
\bfield{Interface} \texttt{firm.tokencomp}: \texttt{Grant(value0, months, cliff, lockup, tax, discount)};
\texttt{schedule(grant)}; \texttt{simulate(grant, vol, drift, n, rng)}; \texttt{summary(sim)}.
\bfield{Rules} Tax is due at vesting on the value then; a tranche is sold at the end of its lock-up; the price model's
drift and volatility are the caller's stated assumptions; the schedule sums to the grant.
\bfield{Acceptance tests} \texttt{code/firm/tokencomp/tests/}: the schedule sums to one with a quarter at the cliff; at
zero volatility the outcome equals the cash value net of tax; volatility spreads the outcome, lowers the median below
the cash value and keeps the mean near it.
\bfield{Stretch} Jumps in the price (Book~4); a tax regime that taxes at sale; hedging a grant with futures or options
where they exist and are allowed.
\end{build}

\begin{omsources}
\item Coinbase Global, Inc., Forms 10-K for 2022, 2023, 2024 and 2025.
\item US Internal Revenue Service, Notice 2014-21; HM Revenue \& Customs, Cryptoassets Manual, CRYPTO21100.
\item Ethereum Foundation, home page and EF Report 2024.
\item US Code, title 11, section 507, and the Judicial Conference's adjustment of dollar amounts effective 1 April 2025.
\end{omsources}

\section{Exercises}

\begin{exercise}[$\star$]\label{exo:in:crypto-native-firms:1}
Compute the exchange's change in revenue and in headcount for each year from 2022 to 2025.
\end{exercise}

\begin{exercise}[$\star$]\label{exo:in:crypto-native-firms:2}
In the example, how much of the grant vests at the cliff, and how much each month after?
\end{exercise}

\begin{exercise}[$\star$]\label{exo:in:crypto-native-firms:3}
An employee receives tokens worth \$10\,000 at vesting, taxed at 40\%, and sells them after a lock-up at half the price.
What is the net outcome?
\end{exercise}

\begin{exercise}[$\star\star$]\label{exo:in:crypto-native-firms:4}
Why is the median outcome of the grant below the cash value while the mean is close to it?
\end{exercise}

\begin{exercise}[$\star\star$]\label{exo:in:crypto-native-firms:5}
A failed firm owes an employee \$60\,000 of wages earned in the last four months. How much has priority under the US rule?
\end{exercise}

\begin{exercise}[$\star\star$]\label{exo:in:crypto-native-firms:6}
From \cref{fig:in:crypto-native-firms:cycle}, in which two years did headcount change most out of step with
revenue, and what explains both?
\end{exercise}

\begin{exercise}[$\star\star\star$]\label{exo:in:crypto-native-firms:7}
\emph{Coding.} Rerun the grant with a twelve-month lock-up. How do the median and the share of paths in which a
tranche's tax exceeds its proceeds change?
\end{exercise}

\begin{exercise}[$\star\star\star$]\label{exo:in:crypto-native-firms:8}
\emph{Find the flaw.} ``The token grant's expected value equals the cash offer, so the two offers are equivalent.''
\end{exercise}

\section{Problem: Paid in Tokens}

\begin{problem}[{Weekend problem --- paid in tokens}]\label{pb:in:crypto-native-firms:1}
An engineer is offered a salary plus a four-year token grant worth \$400\,000 at today's price, or the same salary plus
\$400\,000 in cash over four years. She wants to understand the difference.

\medskip\noindent\textbf{Part I --- The employers.}
\begin{enumerate}
\item Name the four kinds of crypto-native employer.
\item Define a \omterm{def:in:crypto-native-firms:foundation}{protocol foundation} and describe how one foundation funds itself through the cycle.
\item Give four ways crypto-native employers differ from the rest of the industry.
\item Give the exchange's revenue and headcount from 2021 to 2025.
\item What were the two restructurings, in share and number of employees?
\end{enumerate}

\medskip\noindent\textbf{Part II --- Tax.}
\begin{enumerate}[resume]
\item Define \omterm{def:in:crypto-native-firms:tokens}{token compensation} and a \omterm{def:in:crypto-native-firms:tokens}{token warrant}.
\item How do the US and UK tax authorities treat tokens received for work?
\item When is the tax due, and on what value?
\item What happens if the token falls during the lock-up?
\item Compute the net outcome of a \$10\,000 tranche taxed at 40\% and sold at half its vesting price.
\end{enumerate}

\medskip\noindent\textbf{Part III --- The grant.}
\begin{enumerate}[resume]
\item Write the vesting schedule.
\item State the method for valuing the grant.
\item Give the median, the 10th and 90th percentiles and the mean at 80\% volatility.
\item What is the cash equivalent, and in what share of paths is the grant worth less?
\item In what share of paths does at least one tranche's tax exceed its proceeds, at 80\% and at 40\% volatility?
\end{enumerate}

\medskip\noindent\textbf{Part IV --- The verdict.}
\begin{enumerate}[resume]
\item State the \emph{named result}: the median and the 10th--90th percentile range of the grant at 80\% volatility,
against cash, and the share of paths in which tax exceeds proceeds.
\item How much of her unpaid wages would have priority if the firm failed?
\item What does the headcount history suggest about when to join?
\item Name two ways to reduce the grant's risk, if allowed.
\item In two sentences, how should she choose?
\end{enumerate}
\end{problem}

\section{Interview questions}

\begin{interviewq}[$\star$ \iqroles{developer}]\label{iq:in:crypto-native-firms:1}
What does a crypto exchange's engineering team do that a stock exchange's does not?
\end{interviewq}

\begin{interviewq}[$\star$ \iqroles{trader}]\label{iq:in:crypto-native-firms:2}
Why do crypto trading firms staff round the clock, and how would you organise a desk to do it?
\end{interviewq}

\begin{interviewq}[$\star\star$ \iqroles{researcher, risk}]\label{iq:in:crypto-native-firms:3}
A token's price is lognormal with 80\% volatility and no drift. What is the median of its price in one year, relative to
today?
\end{interviewq}

\begin{interviewq}[$\star\star$ \iqroles{risk}]\label{iq:in:crypto-native-firms:4}
You are paid in a token you cannot sell for a year. How could you hedge, and what would stop you?
\end{interviewq}

\begin{interviewq}[$\star\star$ \iqroles{developer, risk}]\label{iq:in:crypto-native-firms:5}
A firm keeps most client assets in cold wallets and some in hot wallets. What decides the split, and what can go wrong?
\end{interviewq}

\begin{interviewq}[$\star\star\star$ \iqroles{researcher}]\label{iq:in:crypto-native-firms:6}
With five years of one firm's revenue and headcount, how would you test whether headcount lags revenue, and what would
you conclude?
\end{interviewq}
